\emph{Two facts used throughout.} (a) For a basis vector $e_{a}$ and a basis blade $e_{B}$ (index set $B$), $e_{a}e_{B} = (-1)^{|B| - [a \in B]}e_{B}e_{a}$: $e_{a}$ anticommutes with each factor of $e_{B}$ other than itself (Theorem \claimref{std-algebra.orthogonal-anticommute}) and commutes with itself. (b) Conjugation by a basis blade maps every basis blade to plus or minus itself: by the product formula of Theorem \claimref{std-algebra.orthogonal-anticommute}, $e_{E}e_{B} = \pm e_{B}e_{E}$, both products being a sign times $(\prod_{E \cap B}Q)\,e_{E \triangle B}$. Hence if $x = \sum_{B}x_{B}e_{B}$ commutes with $e_{E}$, then $0 = e_{E}x - xe_{E} = \sum_{B}x_{B}(\epsilon_{B} - 1)e_{B}e_{E}$ with $\epsilon_{B} = \pm 1$, and since the $e_{B}e_{E}$ are, up to nonzero factors, distinct basis blades, $x_{B} = 0$ whenever $\epsilon_{B} = -1$. So the commutant of any set of basis blades is spanned by the basis blades that commute with all of them, and the commutant of a set of generators can be computed blade by blade. Every commutant below is the commutant of the algebra generated by a set of basis blades, so (b) applies to each.

\emph{(i).} In $\Cl(N,N)$, $n = 2N \ge 2$, the even subalgebra is spanned by even blades, products of the bivectors $e_{a}e_{b}$, so its commutant is the commutant of the bivectors. By (a), $e_{a}e_{b}e_{B} = (-1)^{2|B| - [a \in B] - [b \in B]}e_{B}e_{a}e_{b}$, so $e_{B}$ commutes with $e_{a}e_{b}$ exactly when $[a \in B] = [b \in B]$. Requiring this for every pair $a \ne b$ forces $B = \emptyset$ or $B = \{1, \dots, n\}$. In the order $e_{1}f_{1}e_{2}f_{2}\cdots$ the full blade is $\prod_{j}\omega_{j} = \Gamma$, which squares to $+1$ (Theorem \claimref{tower.volume-elements-planes}). So the commutant is $\mathrm{span}\{1, \Gamma\}$; it has no odd part. Its even part is the centre of $\Cl^{+}(N,N)$, which agrees with the parenthesis of Theorem \claimref{synthesis.plane-commutant-even-ground}(ii) at $r = 0$ ($\omega^{2} = (-1)^{r/2} = 1$). For real $\alpha, \beta$, $(\alpha + \beta\Gamma)^{2} = \alpha^{2} + \beta^{2} + 2\alpha\beta\Gamma$; equal to $-1$ would need $\alpha^{2} + \beta^{2} = -1$. The script checks (b) on all pairs of blades and computes the commutant of the bivectors for $N = 1, 2, 3$, and the square at four rational points.

\emph{(ii).} $J^{2} = e_{2}e_{3}e_{2}e_{3} = -e_{2}^{2}e_{3}^{2} = -1$. The $*$ of Theorem \claimref{tower.structure-tower-c-norms}, restricted to the real level, is the antiautomorphism with $x^{*} = x$ for generators of square $+1$ and $x^{*} = -x$ for those of square $-1$; on a blade it is reversion times $(-1)^{\#\{\text{factors of square } -1\}}$. So $J^{*} = e_{3}^{*}e_{2}^{*} = e_{3}e_{2} = -J$, $v^{*} = \tfrac12(e - f) = w$, and the ladder pairs satisfy the CAR (Theorem \claimref{tower.ladder-pairs-satisfy-car-vectors}; directly, $\{w_{j}, v_{k}\} = \tfrac14\{e_{j} - f_{j}, e_{k} + f_{k}\} = \tfrac12(Q(e_{j}) - Q(f_{j}))\delta_{jk} = \delta_{jk}$ and $\{v_{j}, v_{k}\} = \{w_{j}, w_{k}\} = 0$ since $v_{j}, w_{j}$ are null and orthogonal to the other planes). Every plane generator is orthogonal to $e_{2}$ and $e_{3}$, so anticommutes with each and commutes with $J$; $J$ therefore commutes with $B_{N}$. $I_{4}\sigma = e_{0}e_{1}e_{2}e_{3}e_{0}e_{1}$; moving $e_{0}e_{1}$ to the left past $e_{2}e_{3}$ costs four transpositions, so $I_{4}\sigma = \sigma\sigma J = \sigma^{2}J$, and $\sigma^{2} = e_{0}e_{1}e_{0}e_{1} = -e_{0}^{2}e_{1}^{2} = 1$: $J = I_{4}\sigma$, as in Theorem \claimref{synthesis.transverse-bivector-is-the-complex-unit-of-the-radial-atom}(i). The script checks the CAR, $v^{*} = w$, the antiautomorphism property on all blade pairs for $N \le 2$, and the rest for $N = 1, 2, 3$.

\emph{(iii).} Let $P$ be the set of the $2N$ plane indices and $K = \{2, 3\}$. By (b) the commutant of $B_{N}$ (that of its generators) is spanned by the blades $e_{B}$ with, by (a), $|B| - [a \in B]$ even for every $a \in P$. If $B$ contained one plane index and missed another, $|B| - 1$ and $|B|$ would both be even; so $B \cap P = \emptyset$ or $B \supseteq P$. In the first case $B \subseteq K$ with $|B|$ even: $B = \emptyset$ or $K$, giving $1$ and $J$. In the second $|B|$ is odd, $|B| = 2N + |B \cap K|$, so $|B \cap K| = 1$: $B = P \cup \{2\}$ or $P \cup \{3\}$. Since $\Gamma = \pm e_{P}$, these blades are $\pm\Gamma e_{2}$ and $\pm\Gamma e_{3}$. So the commutant is $\mathrm{span}\{1, J, \Gamma e_{2}, \Gamma e_{3}\}$, of dimension $4$. By Theorem \claimref{tower.volume-elements-planes}, $\Gamma^{2} = 1$ and $\Gamma$ commutes with the non-plane generators $e_{2}$, $e_{3}$, so $(\Gamma e_{2})^{2} = \Gamma^{2}e_{2}^{2} = -1$, likewise $(\Gamma e_{3})^{2} = -1$, $(\Gamma e_{2})(\Gamma e_{3}) = e_{2}e_{3} = J$ and $(\Gamma e_{3})(\Gamma e_{2}) = e_{3}e_{2} = -J$: the quaternion relations. This is the twisted copy $A$ of Theorem \claimref{tower.commuting-copies} for ground $\Cl(0,2)$ with the $N$ planes, generated by $e_{2}\Gamma$ and $e_{3}\Gamma$. A general element $q = a + bJ + c\,\Gamma e_{2} + d\,\Gamma e_{3}$ has $q^{2} = a^{2} - b^{2} - c^{2} - d^{2} + 2a(bJ + c\,\Gamma e_{2} + d\,\Gamma e_{3})$, the three units anticommuting pairwise, so $q^{2} = -1$ exactly when $a = 0$ and $b^{2} + c^{2} + d^{2} = 1$. $\Gamma e_{2}$ and $\Gamma e_{3}$ have odd grade $2N + 1$, so the even part is $\mathrm{span}\{1, J\}$ and its units are $\pm J$. For the kernel rotor $R = e^{\varphi J} = \cos\varphi + J\sin\varphi$ (Theorem \claimref{std-algebra.bivector-exponential}), $J$ is fixed, while $\Gamma e_{2}$, $\Gamma e_{3}$ anticommute with $J$ ($Je_{2} = -e_{2}J$, $\Gamma$ commutes with $J$), so $R(\Gamma e_{2})R^{-1} = R^{2}\Gamma e_{2} = \cos 2\varphi\,\Gamma e_{2} + \sin 2\varphi\,J\Gamma e_{2}$ and $J\Gamma e_{2} = \Gamma e_{3}$: a turn through $2\varphi$ of the plane $(\Gamma e_{2}, \Gamma e_{3})$. A unit $bJ + c\,\Gamma e_{2} + d\,\Gamma e_{3}$ is fixed by every such $R$ exactly when $c = d = 0$. The limit: Theorem \claimref{synthesis.plane-commutant-even-ground}(i), for ground signature $(p, q) = (0, 2)$ ($n = 2 \ge 1$), gives the commutant of all planes in the limit as $\Cl^{0}(0,2) = \mathrm{span}\{1, e_{2}e_{3}\} = \mathrm{span}\{1, J\}$. For $N = 1$ the commutant is the centraliser of $e_{0}, e_{1}$ in $\Cl(1,3)$ stored in Theorem \claimref{synthesis.transverse-bivector-is-the-complex-unit-of-the-radial-atom}(i). The contrast: with ground $\Cl(1,3)$ and transport the $N - 1 \ge 1$ tower planes, the same blade argument (now with $P$ the tower-plane indices and the four ground indices free) gives the blades $e_{B}$ with $B \cap P = \emptyset$, $|B|$ even, or $B \supseteq P$, $|B|$ odd: $2^{3} + 2^{3} = 16$ blades, the twisted copy of $\Cl(1,3)$ of Theorem \claimref{tower.commuting-copies}; its even part is $\Cl^{+}(1,3)$, spanned by $1$, the six bivectors $e_{\mu}e_{\nu}$ and $I_{4}$, of which $e_{1}e_{2}$, $e_{1}e_{3}$, $e_{2}e_{3}$ and $I_{4}$ square to $-1$. So $\pm J$ is not the only even unit there, and $I_{4} = J\sigma$ also commutes with every tower plane. The script computes the commutant blades for $N = 1, 2, 3$, the quaternion relations, the rotor action at $(\cos\varphi, \sin\varphi) = (3/5, 4/5)$, the $N = 1$ centraliser, and the contrast at $N = 2$ (16 blades, 8 even, the four even blades of square $-1$).

\emph{(iv).} In the order $e_{0}e_{1}e_{2}e_{3}e_{1}'f_{1}'\cdots$, $I_{n} = \sigma J\prod_{j \ge 1}\omega_{j}$, and $\sigma$, $J$ commute (each of $e_{0}, e_{1}$ commutes with $J$), so $I_{n} = J\sigma\prod_{j \ge 1}\omega_{j} = J\Gamma$, $\sigma$ being the volume element of the radial plane. $J$ commutes with $\Gamma$ (by (ii)), so $I_{n}^{2} = J^{2}\Gamma^{2} = -1$. $\omega^{*} = (ef)^{*} = f^{*}e^{*} = -fe = ef = \omega$, and the $\omega_{j}$ commute pairwise (Theorem \claimref{tower.volume-elements-planes}), so $\Gamma^{*} = \Gamma$ and $I_{n}^{*} = \Gamma^{*}J^{*} = -\Gamma J = -I_{n}$. $n = vw = \tfrac14(e + f)(e - f) = \tfrac14(e^{2} - ef + fe - f^{2}) = \tfrac12(1 - \omega)$, so $\omega = 1 - 2n$ and $\Gamma = \prod_{j}(1 - 2n_{j})$, which is $+1$ on even and $-1$ on odd occupation: the fermion parity. By Theorem \claimref{tower.volume-elements-planes}, $\Gamma$ anticommutes with each plane generator, hence with $v_{j}, w_{j}$, and commutes with products of two. Then $I_{n}\tfrac12(1 \pm \Gamma) = J\tfrac12(\Gamma \pm \Gamma^{2}) = J\tfrac12(\Gamma \pm 1) = \pm J\tfrac12(1 \pm \Gamma)$. The dimension $2N + 2$ is even, so $I_{n}$ anticommutes with every vector (Theorem \claimref{std-algebra.pseudoscalar-commutation}, $k = 1$), and $\operatorname{Ad}(I_{n})$ negates every generator: it is the grade involution. On $B_{N}$, $\operatorname{Ad}(I_{n}) = \operatorname{Ad}(J)\operatorname{Ad}(\Gamma) = \operatorname{Ad}(\Gamma)$, since $J$ commutes with $B_{N}$. For $N = 1$, $\Gamma = \sigma$ and $I_{n} = I_{4}$, and this is $I_{4}P_{\pm} = \pm JP_{\pm}$ of Theorem \claimref{synthesis.transverse-bivector-is-the-complex-unit-of-the-radial-atom}(iii) with $P_{\pm} = \tfrac12(1 \pm \sigma)$. The script checks all of this for $N = 1, 2, 3$.

\emph{(v).} $B_{N}$ is closed under $*$ (the adjoint of a plane generator is $\pm$ itself), and $J$ commutes with $B_{N}$. Then $\phi_{X}^{*} = X^{*} + X = \phi_{X}$ and $\pi_{X}^{*} = (X^{*} - X)^{*}J^{*} = (X - X^{*})(-J) = J(X^{*} - X) = \pi_{X}$. Since $J$ commutes with $X$ and $X^{*}$, $[\phi_{X}, \pi_{X}] = J[X + X^{*}, X^{*} - X] = J([X, X^{*}] - [X^{*}, X]) = 2J[X, X^{*}]$. For $[X, Y] = 0$ also $[X^{*}, Y^{*}] = -[X, Y]^{*} = 0$, so $[\phi_{X}, \pi_{Y}] = J([X, Y^{*}] - [X^{*}, Y]) = J([X, Y^{*}] + [Y, X^{*}])$; $[\phi_{X}, \phi_{Y}] = [X, Y^{*}] + [X^{*}, Y] = [X, Y^{*}] - [Y, X^{*}]$; and $[\pi_{X}, \pi_{Y}] = J^{2}[X^{*} - X, Y^{*} - Y] = [X^{*}, Y] + [X, Y^{*}]$, the same.

\emph{Collective pairs.} With $b_{k} = v_{2k}w_{2k+1}$, $[b_{k}, b_{k}^{*}] = n_{2k} - n_{2k+1}$ and even elements of distinct blocks commute (Theorem \claimref{synthesis.fluctuation-forms-of-the-tower-and-density-wave-bosons-of-the-null-vacuum}(i), (iii)), so $[S, S^{*}] = \sum_{k}(n_{2k} - n_{2k+1}) =: Q$ and $[\phi, \pi] = 2JQ$. From the CAR, $[n_{j}, v_{j}] = v_{j}$, $[n_{j}, w_{j}] = -w_{j}$, and $n_{j}$ commutes with the modes of other planes (it is even and built from plane $j$); so $[Q, b_{k}] = 2b_{k}$ and $[Q, S^{*}] = -2S^{*}$. The $n_{j}$ are commuting idempotents ($n^{2} = v(wv)w = v(1 - vw)w = n$); $n_{2k}P_{\mathrm{pol}} = P_{\mathrm{pol}}$ and $n_{2k+1}P_{\mathrm{pol}} = 0$, so $QP_{\mathrm{pol}} = MP_{\mathrm{pol}}$ and, by $[Q, S^{*}] = -2S^{*}$, $QX_{j} = (M - 2j)X_{j}$, whence $[\phi, \pi]X_{j} = 2(M - 2j)JX_{j}$ ($J$ commutes with $Q$). $w_{2k+1}(1 - n_{2k+1}) = w_{2k+1}w_{2k+1}v_{2k+1} = 0$ gives $SP_{\mathrm{pol}} = 0$. The $b_{k}^{*}$ commute and $(b_{k}^{*})^{2} = 0$, so $(S^{*})^{j} = j!\sum_{|A| = j}\prod_{k \in A}b_{k}^{*}$, which vanishes for $j > M$. For $j \le M$: $P_{\mathrm{pol}} \neq 0$ (it is a product of $N$ commuting factors $\tfrac12(1 \pm \omega_{j})$, with scalar part $2^{-N}$); $b_{k}b_{k}^{*} = v_{2k}w_{2k+1}v_{2k+1}w_{2k} = n_{2k}(1 - n_{2k+1})$ fixes $P_{\mathrm{pol}}$, so $\prod_{k \in A}b_{k}\prod_{k \in A}b_{k}^{*}P_{\mathrm{pol}} = P_{\mathrm{pol}}$ and each term $T_{A} = \prod_{k \in A}b_{k}^{*}P_{\mathrm{pol}}$ is nonzero; and $n_{2k}T_{A} = [k \notin A]\,T_{A}$, so terms with distinct $A$ lie in distinct joint eigenspaces of left multiplication by the commuting $n_{2k}$ and are linearly independent. Hence $X_{j} \neq 0$. This is the hopping ladder of Theorem \claimref{synthesis.fluctuation-forms-of-the-tower-and-density-wave-bosons-of-the-null-vacuum}(ii) with the unit written in. Dividing by $2M$ gives $[\Phi, \Pi]X_{j} = (1 - 2j/M)JX_{j}$. The script checks the identities for $M = 1, 2$, the values $2(M - 2j)$ on $X_{j}$ for $j = 0, \dots, M$, $X_{M+1} = 0$, and the general identity on a generic element of $B_{2}$.

\emph{Density modes.} $[\rho_{q}, \rho_{q}^{*}] = \sum_{k<q}n_{k} - \sum_{k \ge N-q}n_{k}$ is the edge operator of Theorem \claimref{synthesis.fluctuation-forms-of-the-tower-and-density-wave-bosons-of-the-null-vacuum}(v), so $[\phi_{\rho_{q}}, \pi_{\rho_{q}}] = 2J$ times it. For $2q \le N$ the two index sets are disjoint, $E_{q}$ commutes with every $n_{k}$, $n_{k}E_{q} = E_{q}$ for $k < q$ and $n_{k}E_{q} = 0$ for $k \ge N - q$, so the edge operator times $E_{q}$ is $qE_{q}$, and for $\psi = E_{q}\psi$ it gives $q\psi$. $[\rho_{q}, \rho_{q'}] = \sum_{k, l}[v_{k}w_{k+q}, v_{l}w_{l+q'}]$; by the CAR $[v_{k}w_{m}, v_{l}w_{p}] = \delta_{ml}v_{k}w_{p} - \delta_{pk}v_{l}w_{m}$, so the sum is $\sum_{k}v_{k}w_{k+q+q'} - \sum_{l}v_{l}w_{l+q'+q}$, both over the indices with $k + q + q' \le N - 1$ (the intermediate index then lies in the window automatically), and it vanishes. The script checks, on a window of four modes, the two commutators for $q = 1, 2$, their value $2qJ$ on $E_{q}$, $[\rho_{1}, \rho_{2}] = 0$ and the three mixed identities.

\emph{Lemma.} By the identities above with $u$ in place of $J$ (they used only $u^{2} = -1$, $u^{*} = -u$ and that $u$ commutes with $X$ and $Y$), $[\Phi_{q}, \Pi_{q'}] = u([X_{q}, X_{q'}^{*}] + [X_{q'}, X_{q}^{*}])/(2\sqrt{c_{q}c_{q'}}) = u(c_{q} + c_{q'})\delta_{qq'}/(2\sqrt{c_{q}c_{q'}}) = u\delta_{qq'}$, and $[\Phi_{q}, \Phi_{q'}] = [\Pi_{q}, \Pi_{q'}] = (c_{q} - c_{q'})\delta_{qq'}/(2\sqrt{c_{q}c_{q'}}) = 0$. On $\mathcal{D}$, Theorem \claimref{synthesis.fluctuation-forms-of-the-tower-and-density-wave-bosons-of-the-null-vacuum}(vi) gives $[\rho_{q}, \rho_{q'}] = 0$ and $[\rho_{q}, \rho_{q'}^{*}] = q\delta_{qq'}$, which are the hypotheses with $c_{q} = q$.

\emph{(vi).} The map $e_{3} \mapsto -e_{3}$, others fixed, preserves $Q$; on blades put $\alpha(e_{B}) = (-1)^{[3 \in B]}e_{B}$. By the product formula of Theorem \claimref{std-algebra.orthogonal-anticommute}, $e_{I}e_{J'} = \pm(\prod_{I \cap J'}Q)\,e_{I \triangle J'}$ with a sign depending only on the index sets, and $[3 \in I \triangle J'] \equiv [3 \in I] + [3 \in J'] \pmod 2$, so $\alpha(e_{I}e_{J'}) = \alpha(e_{I})\alpha(e_{J'})$: $\alpha$ is an automorphism. Both $\alpha$ and $*$ are diagonal on blades, so they commute and $\alpha$ is a $*$-automorphism. It fixes every plane generator, hence $B_{N}$ pointwise, and fixes $\sigma = e_{0}e_{1}$; $\alpha(J) = e_{2}(-e_{3}) = -J$, $\alpha(I_{4}) = -I_{4}$, $\alpha(I_{n}) = -I_{n}$, and $\alpha(\phi_{X}) = \phi_{X}$, $\alpha(\pi_{X}) = -J(X^{*} - X) = -\pi_{X}$. Any statement made with $B_{N}$, $*$ and the algebra structure is therefore true of $J$ exactly when it is true of $-J$. With $\sigma$ fixed, $J = I_{4}\sigma$ makes the sign of $J$ and the orientation of $I_{4}$ one datum.

The phase circle: the $n_{j}$ are commuting idempotents that commute with $J$, so $(Jn)^{m} = J^{m}n$ for $m \ge 1$ and $e^{\lambda Jn} = 1 + n(e^{\lambda J} - 1) = 1 - n + ne^{\lambda J}$ (the series of Theorem \claimref{std-algebra.bivector-exponential} for an element of square $-1$, here $J$, and $e^{\lambda J} = \cos\lambda + J\sin\lambda$); the factors commute, so $U_{\theta} = e^{\theta JQ/2}$ is the stated product. $U_{\theta}^{*}$ is the same product with $e^{\pm\theta J/2}$ replaced by $e^{\mp\theta J/2}$ ($n^{*} = n$, $J^{*} = -J$), so $U_{\theta}U_{\theta}^{*} = 1$. From $[\tfrac12JQ, S] = JS$, $\operatorname{Ad}(e^{\theta JQ/2})S = e^{\theta\operatorname{ad}(JQ/2)}S = \sum_{m}\theta^{m}J^{m}S/m! = e^{\theta J}S$, and $U_{\theta}S^{*}U_{\theta}^{*} = (e^{\theta J}S)^{*} = S^{*}e^{-\theta J}$. Then $U_{\theta}\phi U_{\theta}^{*} = e^{\theta J}S + S^{*}e^{-\theta J} = \cos\theta\,\phi + \sin\theta\,J(S - S^{*}) = \phi\cos\theta - \pi\sin\theta$, and $U_{\theta}\pi U_{\theta}^{*} = J(S^{*}e^{-\theta J} - e^{\theta J}S) = \cos\theta\,\pi + \sin\theta\,(S^{*} + S) = \phi\sin\theta + \pi\cos\theta$; as conjugation, it keeps $[\phi, \pi]$. $e^{\theta J}$ commutes with $B_{N}$, so $\operatorname{Ad}(e^{\theta J})$ fixes it; on the kernel, $e_{2}$ anticommutes with $J$, so $e^{\theta J}e_{2}e^{-\theta J} = e^{2\theta J}e_{2} = \cos 2\theta\,e_{2} + \sin 2\theta\,e_{3}$ ($Je_{2} = e_{3}$), a turn through $2\theta$, and $e^{\pi J} = -1$ is the rotor of the turn through $2\pi$. The script checks $\alpha$ as a $*$-homomorphism on blades of grade at most $3$ and its values, and, at $M = 2$ and the rational point $(\cos\tfrac{\theta}{2}, \sin\tfrac{\theta}{2}) = (5/13, 12/13)$, the unitarity of $U_{\theta}$, $U_{\theta}SU_{\theta}^{*} = e^{\theta J}S$, the turn of $(\phi, \pi)$, the invariance of $[\phi, \pi]$ and that $\operatorname{Ad}(e^{\theta J})$ fixes every plane generator. $\square$

\emph{Reading (argued).} Premise (N$_{+}$) assumes the orientation of $I_{4}$ relative to time (\claimref{synthesis.n-plus-null-ground-state}); with the radial plane's orientation $\sigma$ fixed, this is the sign of $J$ by (vi). The fermion phase flow on the atom, $v \mapsto e^{-Jt}v$ (\claimref{synthesis.transverse-bivector-is-the-complex-unit-of-the-radial-atom}(iv)), and the boson relation $[\Phi, \Pi] = +J$ of (v) are both written with it, and $\alpha$ reverses both. The identification is a reading of the premise's wording, not a theorem, and this result does not rest on the premise.
